%! AP Statistics — Unit 1, Lesson 1.2 — Exit Ticket KEY
\documentclass[10pt]{article}
\usepackage{apstats-article}
\usepackage{apstats-key}

%=============================================================================
\begin{document}

\pageheader{Unit 1, Lesson 1.2}{Exit Ticket --- Answer Key}
\begin{tabularx}{\linewidth}{X X X}
Name:\;\ans{Key} & Date:\; & Period:\;
\end{tabularx}

\vspace{0.22in}

\begin{tcolorbox}[colback=sky, colframe=navy, boxrule=0.9pt, arc=2mm,
    left=4mm, right=4mm, top=3mm, bottom=3mm]
\small A researcher is studying music listening habits. Classify each variable
and justify.
\end{tcolorbox}

\vspace{0.18in}

\renewcommand{\arraystretch}{2.4}
\begin{tabularx}{\linewidth}{@{} c >{\raggedright\arraybackslash}p{0.36\linewidth}
                                    C{0.08\linewidth} X @{}}
\rowcolor{sky}
\textbf{\#} & \textbf{Variable} & \textbf{Type} & \textbf{Justification} \\
\midrule
1 & Favorite music genre
  & \ans{C}
  & \ans{Labels a category; averaging genres is meaningless} \\
2 & Number of songs on a playlist
  & \ans{QD}
  & \ans{Count of whole songs; gaps exist between values} \\
3 & Duration of the longest song (minutes)
  & \ans{QC}
  & \ans{Measured time; can be any positive decimal (e.g., 4.73 min)} \\
4 & Country of birth of the artist
  & \ans{C}
  & \ans{Places artist into a geographic category; no arithmetic} \\
5 & Temperature at venue at noon (°F)
  & \ans{QC}
  & \ans{Measured on a continuous scale; can take any real value in a range} \\
\end{tabularx}

\vspace{0.2in}

\begin{teachernote}
\textbf{Scoring (completeness):} Award credit for any sincere attempt with a reason.
Full credit requires both a correct type \emph{and} a justification that references
the average test or the countable/measured distinction.\\[4pt]
\textbf{Common errors to flag for re-teaching:}
\begin{itemize}[itemsep=3pt]
  \item \#3 marked QD: Students confuse ``minutes'' with ``whole minutes.'' Clarify that
        song duration can be measured to fractions of a second.
  \item \#2 marked QC: Students think ``number of'' could be fractional. Remind them
        you cannot have 3.7 songs on a playlist.
  \item \#1 or \#4 marked quantitative: Students see words but still hesitate. Ask:
        ``What would it mean to average genres?''
\end{itemize}
\end{teachernote}

\vspace{0.16in}

\begin{tcolorbox}[colback=redbg, colframe=redacc, boxrule=0.8pt, arc=2mm,
    left=4mm, right=4mm, top=3mm, bottom=3mm]
\small\textbf{AP-Style Practice}\\[6pt]
A survey records annual income (dollars) and job type (teacher, engineer, nurse, other).

\vspace{10pt}
\begin{enumerate}[label=\textbf{(\Alph*)}, leftmargin=2em, itemsep=6pt]
  \item Annual income is categorical; job type is quantitative.
  \item Annual income is quantitative discrete; job type is categorical.
  \item \textbf{\textcolor{keyred}{Annual income is quantitative continuous; job type is categorical.}}
        \textcolor{keyred}{\textit{$\leftarrow$ correct}}
  \item Both variables are categorical.
\end{enumerate}
\end{tcolorbox}

\vspace{0.1in}
\noindent\ans{Answer: C} \quad
\ans{Income is a measured amount that can take any non-negative value --- it is quantitative
continuous. Job type places each person into a named category --- it is categorical.
Choice B is wrong because income is not restricted to whole-dollar values in practice;
Choice D incorrectly calls income categorical.}

\vspace{0.12in}
\begin{teachernote}
\textbf{Most common wrong answer: B.} Students who choose B are applying the right framework
(recognizing income is quantitative) but mis-classifying it as discrete because they associate
dollars with whole numbers. Prompt: ``Can income be \$47,382.50?'' This usually resolves it.\\[3pt]
If more than a third of the class chooses D, the categorical/quantitative distinction needs
another pass at the start of Lesson 1.3.
\end{teachernote}

\end{document}
